%% sections/06_conclusions.tex

\section{Conclusions}
\label{sec:conclusions}

We have presented \AAS\ (AstDyn-Adaptive Symplectic), a fourth-order
adaptive geometric integrator for asteroid orbit propagation,
with a symplectic fixed-step DKD kernel and an adaptive
Sundman-type step-size strategy.
The main contributions of this work are:

\begin{enumerate}

  \item \textit{Hessian-based step size function.}
    The step $\Delta t = \eps / \sqrt{\rho(\Hess(\vq))}$
    is physically motivated by local orbital curvature and is computed
    analytically at negligible cost. The induced Keplerian
    $r^{3/2}$ scaling is consistent with prior adaptive-symplectic
    results; the added value here is the Hessian-based formulation and
    its direct integration with analytical STM updates.

  \item \textit{Time-reversibility.}
    Harmonic-mean symmetrization of $g(\vq)$ ensures that the
    discrete map satisfies the time-reversal invariance condition
    of \citet{PretoTremaine1999}.

  \item \textit{Analytical STM propagation.}
    The State Transition Matrix is propagated via Hessian-based
    updates within the Drift-Kick-Drift sequence, eliminating the
    $12N_{\mathrm{steps}}$ overhead of numerical STM approaches;
    this is directly relevant for orbital uncertainty propagation
    \citep{Zhou2024dstt}.

  \item \textit{Shadow Hamiltonian monitoring.}
    The leading $O(\Delta t^4)$ correction to the shadow Hamiltonian
    is computed analytically, providing a rigorous energy health check.

  \item \textit{Numerical validation.}
    Benchmarks on Ceres, Apophis, and Phaethon demonstrate
    (i) one-year \AAS\ energy errors down to
    $1.96\times10^{-13}$ (Apophis),
    (ii) for Phaethon at equal accuracy
    ($|\Delta E/E|\approx7.00315\times10^{-5}$),
    $7.0644\times10^4$ (\AAS, $\varepsilon=10^{-3}$)
    vs.
    $2.922\times10^4$ (\saba) function evaluations
    (a factor $2.42$ more for \AAS, i.e.
    no global speedup claim in this equal-accuracy regime), and
    (iii) analytical-vs-numerical STM agreement
    $\|\STM_{\AAS}-\STM_{\mathrm{num}}\|_F(30\,\mathrm{d})
    =1.45983\times10^{-5}$.

\end{enumerate}

The full paper package, benchmark scripts, and figure data used in this
article are available at \url{https://github.com/ITALOccult/AAS---Paper---MNRAS}.
\AAS\ is implemented within the \astdyn\ library.
A companion paper (Bigi, in prep.) describes the
full library pipeline from MPC observations to stellar occultation
path predictions, validated against JPL Horizons with residuals
$< 2.5\,\mu$m.

For completeness, we note that the present benchmark campaign remains
partly reduced-complexity; full operational validation will require
systematic N-body tests (including Jupiter), strongly perturbed
intervals, and close-encounter cases.
For the STM contribution specifically, the current evidence is a
single-case consistency test; a full validation still requires
step-size convergence, numerical-STM cross-comparisons, and
orbit-determination impact studies.
Likewise, dedicated Kahan-ablation experiments and formal statistical
uncertainty/sensitivity analyses remain necessary before elevating
the present deterministic benchmark outcomes to fully statistical
performance statements.
